\noindent Patch tokenisation is the input stage of time-series foundation models. When a signal completes a whole number of cycles within one stride or one patch, the token sequence repeats. The frequencies at which this happens, called candidate frequencies here, follow from the sampling rate $f_s$, the patch size $P$ and the stride $S$ alone, and mark where the model could become blind. This report asks whether Chronos-Bolt-Tiny, retrained under fifteen geometries, loses information at those frequencies, in a photovoltaic control setting assessed under the AI Risk Management Framework of the National Institute of Standards and Technology.

\noindent Three hypotheses are tested: that blind spots exist, in its forecast or its internal state; that such a loss does not depend on the phase of the signal; and that the candidate frequencies move with the geometry. Overlapping patches is also tested as a mitigation.

\noindent Three exploratory analyses came first. A sweep of pure tones through every geometry found that the forecast rebuilds a tone at the right frequency only in the lower part of the band, while linear probes on the internal state still identify most frequencies, and that the forecast is drawn towards candidate frequencies rather than losing them. Single tones added to generated signals were recovered in about a quarter of the trials at a candidate frequency and almost never beside one. On a year of measured photovoltaic power no retrained geometry beat the published $P=S=16$, and the dominant daily cycle lies far below the first candidate frequency.

\noindent Before the pre-registered gates, the Bayesian posteriors also point away from a blind spot: at a candidate frequency the forecast recovers a tone $35\%$ better than beside it, and the internal state needs only $6\%$ more bits to read it. The effect barely changes with phase, the candidate frequencies tied to the stride move exactly as predicted, those tied to the patch could not be isolated, and overlap brings no measurable mitigation. All fits converged, but none reproduces the dispersion of the data and three models also fail parameter recovery, so no claim is reportable.

\noindent The risk assessment therefore keeps open every exposure that depends on a loss and monitors the candidate frequencies on the raw signal, before tokenisation.
